\clearpage
\chapter*{Abstract}
\addcontentsline{toc}{chapter}{Abstract}

Concrete spalling, the sudden explosive detachment of concrete surface when concrete is exposed to rapidly rising temperature. This phenomenon is a major threat to the integrity of structure subjected to fire or to severe accident conditions, such as nuclear containtment buildings. It results from coupled thermo-hydro-mechanical (THM) process of the material, which is staill incompletely understood and remain difficult to predict.

The work studies the THM behaviors of concrete under severe accident conditions using a stochastic poro-mechanical approach. A numerical model of a coupled THM model is implemented in the finite element code \textsc{Cast3M}. In which the thermo-hydric subsystem is solved monolithically and the mechanical problem (linear elasticity with a \textsc{Mazars} damage law) is coupled in a staggered manner.

Unlike the traditional approach of TH coupling in \cite{dauti2018}, this work adopts the new thermo-hydro-chemical (THC) approach of \cite{sciume2024}. The transport and storage properties are redefined using the \emph{hydration degree}, this ensures that the initial hydro-chemical state of the material is intergrated into the model.

The model is then calibrated and validated against the experiment conducted by Kalifa \cite{kalifa2000}, used as a reference benchmark, with the aim of reproducing the measured temperature and gas pressure evolutions within concrete specimens. The temperature field is reproduced accurately, whereas the gas pressure cannot be matched with the same confidence. This raise the question of whether the gauges truly recorded the pore gas pressure.

Finally, the outcomes of the calibrated model contribute as the input for a stochastic assessment of the model, carried out with \textsc{OpenTURNS}. Therefore, it provides insigts of the correlation and sensitivity analyses identify the dominant parameters, a reliability analysis quantifies the probability of failure and the risk of spalling, and a random-field description of material heterogeneity is shown to produce more critical predictions than the homogeneous assumption.

